% ======================================================================
%  CHƯƠNG 5 — KẾT LUẬN VÀ HƯỚNG PHÁT TRIỂN TƯƠNG LAI
% ======================================================================
\chapter{Kết luận và hướng phát triển tương lai}
\label{ch:ketluan}

\section{Kết luận}
\label{sec:ketluan}
Đồ án đã xây dựng và vận hành một quy trình đánh giá công bằng cho các thư viện
học máy tự động, đồng thời áp dụng quy trình đó để so sánh bốn thư viện AutoML mã
nguồn mở trên 12 bộ dữ liệu thực thu thập từ Kaggle.

Về mặt phương pháp luận, đồ án đã hệ thống hoá bốn điều kiện của một phép so sánh
công bằng và hiện thực hoá từng điều kiện bằng một cơ chế kiểm soát cụ thể: phân
hoạch fold được sinh một lần từ hạt giống cố định và dùng chung cho mọi thư viện;
mỗi thư viện được cô lập trong một môi trường ảo riêng nhằm loại trừ xung đột phụ
thuộc; ngân sách thời gian được áp đặt đồng nhất ở cả hai lớp thư viện và tiến
trình; và mọi lần chạy thất bại đều được ghi nhận kèm phân loại nguyên nhân thay
vì bị loại khỏi thống kê.

Về mặt hạ tầng, nhóm đã xây dựng bộ công cụ điều phối theo mô hình tiến trình chủ
--- tiến trình con, hỗ trợ phân tán khối lượng tính toán trên nhiều máy, khôi phục
sau gián đoạn và hợp nhất kết quả từng phần. Hạ tầng này đã hoàn tất khối lượng
thực nghiệm gồm bốn thư viện trên 12 bộ dữ liệu với 5 fold và hai mức ngân sách
thời gian.

Về mặt kết quả, phát hiện quan trọng nhất là \emph{không tồn tại thư viện AutoML
vượt trội toàn diện}. Mô hình cây Bradley--Terry cho thấy thứ hạng đảo chiều rõ
rệt giữa các dạng tác vụ: H2O~AutoML dẫn đầu ở phân loại đa lớp và hồi quy nhưng
xếp cuối ở phân loại nhị phân, trong khi MLJAR đứng đầu nhóm nhị phân nhưng xếp
cuối nhóm đa lớp. Một xếp hạng toàn cục đơn lẻ, nếu được sử dụng, sẽ che khuất
hoàn toàn cấu trúc này.

Bên cạnh đó, thực nghiệm cho thấy chênh lệch về chất lượng dự đoán giữa các thư
viện thường nhỏ hơn độ lệch chuẩn giữa các fold, trong khi chênh lệch về chi phí
tài nguyên lại rất lớn: thời gian huấn luyện chênh nhau tới bốn lần và bộ nhớ tới
hơn mười lần cho cùng một mức chất lượng. Điều này gợi ý rằng trong thực tiễn
triển khai, các tiêu chí về chi phí, độ ổn định và tỉ lệ thất bại nên được cân
nhắc ngang bằng, nếu không muốn nói là trước, so với chỉ số chất lượng thuần tuý.

\section{Trả lời các câu hỏi nghiên cứu}
\label{sec:traloi}
Đối chiếu với mục tiêu đặt ra ở Chương~\ref{ch:gioithieu}, các kết quả thu được
cho phép trả lời ba câu hỏi trung tâm của đồ án.

\paragraph{Có tồn tại một thư viện AutoML tốt nhất hay không?} Câu trả lời là
không. Mô hình cây Bradley--Terry cho thấy H2O~AutoML đạt cường độ $0.446$ ở nhóm
phân loại đa lớp và $0.389$ ở nhóm hồi quy --- dẫn đầu cả hai --- nhưng chỉ đạt
$0.164$ ở nhóm phân loại nhị phân, tức xếp cuối. MLJAR thể hiện xu hướng ngược
lại. Nếu chỉ nhìn vào bảng xếp hạng toàn cục, nơi ba thư viện dẫn đầu chênh nhau
chưa tới $0.03$, cấu trúc đảo chiều này sẽ hoàn toàn bị che khuất. Kết luận rút ra
là việc lựa chọn thư viện phải xuất phát từ dạng tác vụ cụ thể, và các bảng xếp
hạng AutoML rút gọn thành một con số duy nhất cần được tiếp nhận thận trọng.

\paragraph{Chênh lệch giữa các thư viện nằm ở đâu?} Không nằm chủ yếu ở chất lượng
dự đoán. Trên nhiều bộ dữ liệu, khác biệt điểm số giữa thư viện tốt nhất và kém
nhất nhỏ hơn độ lệch chuẩn giữa các fold của chính từng thư viện --- trường hợp
\texttt{adult\_income} với chênh lệch AUC $0.0018$ so với độ lệch chuẩn
$0.0015$--$0.0019$ là ví dụ điển hình. Trong khi đó, chênh lệch về thời gian huấn
luyện lên tới bốn lần và về bộ nhớ lên tới hơn mười lần. Do vậy, trong phần lớn
tình huống ứng dụng, tiêu chí chi phí và độ ổn định có sức phân biệt cao hơn tiêu
chí chất lượng.

\paragraph{Cấp thêm tài nguyên có luôn cải thiện kết quả không?} Không phải lúc
nào cũng vậy. Khi tăng ngân sách từ $60$ lên $300$ giây, MLJAR cải thiện mạnh về
chất lượng (từ $0.59$ lên $0.79$ trên thang chuẩn hoá) nhưng đồng thời tăng tỉ lệ
thất bại từ $0\%$ lên khoảng $20\%$. Việc nới lỏng ràng buộc thời gian đã kích
hoạt các chế độ tìm kiếm tốn bộ nhớ hơn, qua đó chuyển nút thắt cổ chai từ thời
gian sang bộ nhớ. Quan hệ giữa tài nguyên và hiệu năng, vì vậy, không đơn điệu và
cần được kiểm soát đồng thời trên nhiều chiều.

\section{Giới hạn của nghiên cứu}
\label{sec:hanche}
Nhóm ghi nhận một số giới hạn của kết quả trình bày trong báo cáo.

\textbf{Về phép đo bộ nhớ.} Mức tiêu thụ bộ nhớ được quan sát từ tiến trình
Python, do đó không phản ánh đúng thực tế đối với H2O~AutoML --- thư viện thực
hiện huấn luyện trong một tiến trình máy ảo Java riêng biệt. Vì lý do này, đồ án
không xếp hạng các thư viện theo tiêu chí bộ nhớ, và các con số của H2O ở
Mục~\ref{sec:bonho} chỉ nên được hiểu là chi phí phía máy khách.

\textbf{Về tuyến phương pháp cơ sở.} Đây là giới hạn đáng kể nhất của nghiên cứu.
Trong tổng số 120 lần chạy đã lên lịch cho hai phương pháp cơ sở, chỉ 15 lần hoàn
tất; phần còn lại hoặc dừng giữa chừng, hoặc chưa được thực thi do quỹ thời gian
tính toán mà nhóm huy động được đã dùng gần hết cho bốn thư viện AutoML chính.
Riêng rừng ngẫu nhiên còn vướng một khiếm khuyết ở lớp bao phương pháp cơ sở như
đã nêu ở Mục~\ref{sec:baseline}. Hệ quả là đối chiếu ở Bảng~\ref{tab:baseline}
mới thực hiện được với một phương pháp cơ sở trên ba bộ dữ liệu, thay vì hai
phương pháp trên toàn bộ mười hai bộ như thiết kế ban đầu.

\textbf{Về quy mô ngân sách.} Thực nghiệm giới hạn ở hai mức ngân sách $60$ và
$300$ giây. Các nghiên cứu quy mô lớn thường sử dụng ngân sách hàng giờ, nơi trật
tự giữa các thư viện có thể thay đổi --- đặc biệt với những thư viện còn nhiều dư
địa cải thiện như MLJAR.

\textbf{Về mức độ cô lập môi trường.} Bộ kết quả trong báo cáo được tạo bằng đường
dòng lệnh, vốn cô lập các thư viện ở mức môi trường ảo Python. Lựa chọn này cho
phép chạy phân tán trên nhiều máy nhưng không cô lập tầng thư viện hệ thống như
container. Điều kiện công bằng vẫn được bảo toàn trong mỗi lượt chạy vì mọi thư
viện dùng chung một cấu hình phần cứng; đổi lại, việc tái lập chính xác trên máy
khác đòi hỏi thêm thao tác thiết lập. Đường chạy trên nền Docker của AMLB Studio
khắc phục hạn chế này nhưng chưa được dùng cho lượt đánh giá hiện tại.

\textbf{Về ý nghĩa thống kê.} Các chênh lệch nhỏ giữa những thư viện dẫn đầu chưa
được kiểm chứng bằng kiểm định hình thức. Báo cáo đã thận trọng không tuyên bố
thắng thua khi chênh lệch nằm dưới ngưỡng độ lệch chuẩn giữa các fold, nhưng đây
là biện pháp phòng vệ chứ không thay thế được một kiểm định đúng nghĩa.

\textbf{Về phạm vi bài toán.} Nghiên cứu giới hạn ở dữ liệu dạng bảng với ba dạng
tác vụ học có giám sát; các kết luận không mở rộng trực tiếp sang dữ liệu ảnh, văn
bản hay chuỗi thời gian.

\section{Hướng phát triển tương lai}
\label{sec:huongphattrien}

\subsection{Hoàn tất tuyến phương pháp cơ sở}
Như đã ghi nhận ở Mục~\ref{sec:baseline}, tuyến phương pháp cơ sở chưa chạy trọn
vẹn do quỹ thời gian tính toán của nhóm đã dành phần lớn cho bốn thư viện AutoML
chính, và riêng rừng ngẫu nhiên còn vướng lỗi mã hoá dữ liệu đầu vào. Khắc phục
không đòi hỏi thay đổi giao thức mà chỉ cần hai việc: bổ sung một bước tiền xử lý
tối thiểu và đồng nhất --- mã hoá one-hot cho biến hạng mục, điền giá trị khuyết
bằng trung vị --- vào lớp bao phương pháp cơ sở, rồi chạy lại trên toàn bộ 12 bộ
dữ liệu. Chi phí tính toán cho bước này thấp hơn nhiều so với một lượt chạy AutoML,
bởi cả hai phương pháp cơ sở đều không thực hiện tìm kiếm siêu tham số.

Giá trị của việc này vượt xa một dòng số liệu bổ sung. Rừng ngẫu nhiên với cấu
hình mặc định là ngưỡng so sánh có ý nghĩa thực tiễn nhất trong toàn bộ nghiên
cứu: nó lượng hoá phần lợi ích mà quá trình tìm kiếm tự động mang lại so với
phương án đơn giản nhất mà một người thực hành có kinh nghiệm sẽ chọn. Nếu khoảng
cách đó hẹp trên phần lớn bộ dữ liệu, kết luận thực tiễn của đề tài sẽ thay đổi
đáng kể --- và đó là câu hỏi mà bộ kết quả hiện tại chưa trả lời được.

\subsection{Áp dụng kiểm định thống kê lên bộ kết quả hiện có}
Quy trình kiểm định Friedman và hậu kiểm Nemenyi đã được hiện thực trong lớp phân
tích của hệ thống (Mục~\ref{sec:console}) nhưng chưa được áp dụng lên bộ kết quả ở
Chương~\ref{ch:thucnghiem}, do hai phần được phát triển song song trên hai nhánh
khác nhau. Kết nối chúng là bước tiếp theo trực tiếp nhất và không đòi hỏi chạy
lại thực nghiệm.

Với bốn thư viện trên 12 bộ dữ liệu, sau khi thu về tập con đầy đủ vẫn còn chín bộ,
vượt ngưỡng tối thiểu mà kiểm định yêu cầu. Kết quả sẽ cho phép phát biểu dứt khoát
hai điều mà báo cáo hiện chỉ nhận định dựa trên độ lệch chuẩn quan sát: cách biệt
giữa H2O và FLAML ở nhóm hồi quy có phải khác biệt thực hay không, và khoảng cách
$0.004$ giữa H2O và MLJAR ở mức toàn cục có nằm dưới ngưỡng phân biệt hay không.
Biểu đồ khoảng cách tới hạn đi kèm sẽ thay cách trình bày thứ hạng hiện tại bằng
một hình thức chặt chẽ hơn, trong đó những thư viện không phân biệt được về mặt
thống kê được nối với nhau thay vì bị xếp thứ tự dứt khoát.

\subsection{Mở rộng thang ngân sách và khảo sát đường cong hội tụ}
Thay vì chỉ so sánh hai mức ngân sách rời rạc, hướng phát triển tiếp theo là ghi
nhận đường cong hiệu năng theo thời gian trong suốt quá trình tìm kiếm của mỗi thư
viện. Cách tiếp cận này cho phép trả lời những câu hỏi mà thiết kế hiện tại chưa
xử lý được: mỗi thư viện cần bao nhiêu thời gian để đạt $95\%$ hiệu năng cuối
cùng, và thời điểm nào hiệu suất biên của việc cấp thêm tài nguyên trở nên không
đáng kể. Với kết quả cho thấy MLJAR cải thiện mạnh từ $0.59$ lên $0.79$ khi tăng
ngân sách, việc xác định điểm bão hoà của thư viện này có giá trị thực tiễn trực
tiếp trong việc hoạch định tài nguyên.

\subsection{Kiểm soát đồng thời nhiều chiều tài nguyên}
Phát hiện về việc MLJAR tăng tỉ lệ thất bại khi được cấp thêm thời gian cho thấy
ngân sách thời gian và ngân sách bộ nhớ tương tác với nhau theo cách không hiển
nhiên. Hướng nghiên cứu tiếp theo là thiết kế thực nghiệm hai chiều, trong đó cả
giới hạn thời gian lẫn giới hạn bộ nhớ đều được biến thiên có hệ thống, nhằm xây
dựng bản đồ vùng vận hành an toàn cho từng thư viện. Kết quả sẽ trả lời câu hỏi
mà người triển khai thực sự quan tâm: với một cấu hình máy cho trước, mức ngân
sách thời gian nào là tối ưu mà vẫn giữ tỉ lệ thất bại ở mức chấp nhận được.

\subsection{Đo bộ nhớ ở mức tiến trình hệ điều hành}
Để khắc phục giới hạn của phép đo hiện tại, mức tiêu thụ bộ nhớ nên được ghi nhận
ở mức tiến trình hệ điều hành thay vì từ bên trong tiến trình Python, qua đó bao
gồm cả bộ nhớ do máy ảo Java của H2O cấp phát. Khi đó, bộ nhớ mới có thể trở thành
một tiêu chí xếp hạng hợp lệ bên cạnh chất lượng và thời gian, và bức tranh đánh
đổi giữa các thư viện sẽ hoàn chỉnh hơn.

\subsection{Tích hợp meta-learning để khuyến nghị tự động}
Kết quả của đồ án cho thấy thư viện phù hợp phụ thuộc vào dạng tác vụ và đặc điểm
dữ liệu. Điều này gợi mở khả năng xây dựng một cơ chế khuyến nghị tự động: từ các
đặc trưng mô tả bộ dữ liệu (số mẫu, số đặc trưng, tỉ lệ biến hạng mục, mức mất cân
bằng nhãn, dạng tác vụ), huấn luyện một mô hình dự đoán thư viện nào có khả năng
cho kết quả tốt nhất trong ngân sách cho trước. Bộ kết quả gồm 12 bộ dữ liệu hiện
tại còn quá nhỏ để huấn luyện một mô hình như vậy, song việc mở rộng tập dữ liệu
lên vài chục bộ sẽ tạo đủ cơ sở. Hướng phát triển này biến kết quả đánh giá từ một
bảng tra cứu tĩnh thành một công cụ hỗ trợ ra quyết định.

\subsection{Mở rộng tập thư viện và dạng dữ liệu}
Cuối cùng, việc bổ sung các thư viện theo hướng tiếp cận khác biệt --- chẳng hạn
auto-sklearn với tối ưu Bayes kết hợp meta-learning, hay TPOT với giải thuật tiến
hoá --- sẽ giúp kiểm chứng liệu hiện tượng đảo chiều xếp hạng theo dạng tác vụ có
phải là đặc tính chung của lĩnh vực AutoML hay chỉ giới hạn trong bốn thư viện đã
khảo sát. Song song, việc mở rộng sang dữ liệu chuỗi thời gian sẽ cho phép đánh
giá năng lực của các thư viện trong một lớp bài toán có cấu trúc phụ thuộc thời
gian, nơi giả định về tính độc lập giữa các mẫu không còn đúng và do đó bản thân
giao thức kiểm chứng chéo cũng cần được thiết kế lại.
